\clearpage
\chapter*{Símbolos y abreviaturas}
\addcontentsline{toc}{chapter}{Símbolos y abreviaturas}
\markboth{Símbolos y abreviaturas}{}

\section*{Símbolos}

\small
\renewcommand{\arraystretch}{0.88}
\begin{longtable}{@{}p{0.20\textwidth}p{0.74\textwidth}@{}}
\toprule
Símbolo & Significado \\
\midrule
\endhead
\(\bm{B}\)                 & Vector campo magnético \\
\(|\vec{B}|\)              & Módulo del campo magnético \\
\(B_z\)                    & Componente del campo normal a la cara del imán (magnitud medida por la sonda Hall) \\
\(B_r\)                    & Remanencia del imán; único parámetro libre del modelo numérico \\
\(\partial B/\partial z\)  & Gradiente vertical del campo en el centro de la capa de semillas \\
\(\mu_{\mathrm B}\)        & Magnetón de Bohr, \(9{,}27\times10^{-24}\,\mathrm{J\,T^{-1}}\) \\
\(k_{\mathrm B}\)          & Constante de Boltzmann \\
\(T\)                      & Temperatura absoluta \\
\(\Delta\)                 & Imbibición: cambio relativo de peso entre el estado seco y las 72~h; variable de respuesta primaria (Ec.~\ref{eq:pct_change}) \\
\(w_{\text{inicial}},\,w_{\text{final}}\) & Masa del lote de semillas en seco y a las 72\,h \\
\(n\)                      & Número de unidades experimentales (muestras) \\
\(F_{\nu_1,\nu_2}\)        & Estadístico \(F\) con \(\nu_1\) y \(\nu_2\) grados de libertad \\
\(p\)                      & Valor \(p\) \\
\(\alpha\)                 & Nivel de significación (0,05) \\
\(\sigma\)                 & Desviación estándar \\
\(R^{2}\)                  & Coeficiente de determinación \\
\bottomrule
\end{longtable}

\vspace{0.4em}

\section*{Abreviaturas}

\begin{longtable}{@{}p{0.20\textwidth}p{0.74\textwidth}@{}}
\toprule
Abreviatura & Significado \\
\midrule
\endhead
ANOVA   & Análisis de la varianza (\textit{analysis of variance}) \\
CME     & Campo magnético estático \\
FAD     & Flavina adenina dinucleótido \\
FDR     & Tasa de falsos descubrimientos (\textit{false discovery rate}) \\
HSD     & Diferencia honestamente significativa (\textit{honestly significant difference}) \\
MBLab   & Laboratorio de Magnetobiología, UNSL \\
RMSE    & Raíz del error cuadrático medio (\textit{root mean square error}) \\
ROS     & Especies reactivas del oxígeno (\textit{reactive oxygen species}) \\
RPM     & Mecanismo de pares de radicales (\textit{radical pair mechanism}) \\
SD      & Desviación estándar (\textit{standard deviation}) \\
UNSL    & Universidad Nacional de San Luis \\
\bottomrule
\end{longtable}
\renewcommand{\arraystretch}{1}
\normalsize

\endinput
